\chapter{图论}\label{ux56feux8bba}

\section{图的表示与基础遍历}\label{ux56feux7684ux8868ux793aux4e0eux57faux7840ux904dux5386}

稀疏图用邻接表，边结构保存 \passthrough{\lstinline!to, weight, id!}；需要删除/反向边时给每条有向边保存反向边下标。无向边应添加两条边。网格可把 \passthrough{\lstinline!(x,y)!} 映射为 \passthrough{\lstinline!x*m+y!}，或直接四方向 BFS。

DFS 可求连通块、染色、子树大小、DFS 序、树高；BFS 求无权最短路和层数。递归 DFS 在 \passthrough{\lstinline!n≈2e5!} 时可能爆栈，改手写栈或提高栈空间。

\begin{lstlisting}[language={C++}]
vector<vector<int>> g(n);
vector<int> dist(n, -1);
queue<int> q;
dist[s] = 0; q.push(s);
while (!q.empty()) {
    int u = q.front(); q.pop();
    for (int v : g[u]) if (dist[v] == -1)
        dist[v] = dist[u] + 1, q.push(v);
}
\end{lstlisting}

\section{有向无环图（Directed Acyclic Graph，DAG）与拓扑排序}\label{ux6709ux5411ux65e0ux73afux56fedirected-acyclic-graphdagux4e0eux62d3ux6251ux6392ux5e8f}

Kahn：把入度为 0 的点入队，删除出边；处理点数小于 \passthrough{\lstinline!V!} 说明有环。DFS 三色标记也能判环。DAG 上按拓扑序做最长路/最短路、路径计数、DP；多源最短路可把所有源同时入队。

关键路径：\passthrough{\lstinline!earliest[v] = max(earliest[v], earliest[u]+w)!}。若要恢复方案，保存前驱；同一最优值的多条路径计数要注意模数。

\begin{lstlisting}[language={C++}]
queue<int> q;
for (int u = 0; u < n; ++u) if (!indeg[u]) q.push(u);
vector<int> topo;
while (!q.empty()) {
    int u = q.front(); q.pop(); topo.push_back(u);
    for (auto [v, w] : g[u]) if (--indeg[v] == 0) q.push(v);
}
if ((int)topo.size() != n) throw runtime_error("cycle");
\end{lstlisting}

\section{最短路}\label{ux6700ux77edux8def}

\begin{longtable}[]{@{}lll@{}}
\toprule\noalign{}
条件 & 算法 & 复杂度 \\
\midrule\noalign{}
\endhead
\bottomrule\noalign{}
\endlastfoot
无权 & BFS & \passthrough{\lstinline!O(V+E)!} \\
边权 0/1 & 0-1 BFS & \passthrough{\lstinline!O(V+E)!} \\
非负权 & Dijkstra + 堆 & \passthrough{\lstinline!O((V+E)log V)!} \\
允许负边、单源 & Bellman-Ford & \passthrough{\lstinline!O(VE)!} \\
DAG & 拓扑 DP & \passthrough{\lstinline!O(V+E)!} \\
全源、\passthrough{\lstinline!V≤500!} & Floyd & \passthrough{\lstinline!O(V³)!} \\
\end{longtable}

Dijkstra 的核心前提是所有边权非负；堆中允许重复状态，弹出时若距离不是当前 \passthrough{\lstinline!dist[u]!} 就跳过。多源时将所有起点距离设为 0。

负环检测：Bellman-Ford 第 \passthrough{\lstinline!V!} 轮仍有更新即存在从源可达负环；SPFA 可做队列优化但最坏仍为 \passthrough{\lstinline!O(VE)!}，不要无条件依赖。

差分约束：不等式 \passthrough{\lstinline!x\_v ≤ x\_u + w!} 建 \passthrough{\lstinline!u→v!} 权 \passthrough{\lstinline!w!}，求最短路；若存在负环则无解。求最大下界时改写方向/用最长路。

\begin{lstlisting}[language={C++}]
using P = pair<long long, int>;
vector<long long> dist(n, INF);
priority_queue<P, vector<P>, greater<P>> pq;
dist[s] = 0; pq.push({0, s});
while (!pq.empty()) {
    auto [du, u] = pq.top(); pq.pop();
    if (du != dist[u]) continue;
    for (auto [v, w] : g[u]) if (dist[v] > du + w) {
        dist[v] = du + w;
        pq.push({dist[v], v});
    }
}
\end{lstlisting}

\begin{lstlisting}[language={C++}]
deque<int> q;
vector<int> d(n, INF); d[s] = 0; q.push_front(s);
while (!q.empty()) {
    int u = q.front(); q.pop_front();
    for (auto [v, w] : zero_one_graph[u]) if (d[v] > d[u] + w) {
        d[v] = d[u] + w;
        if (w == 0) q.push_front(v); else q.push_back(v);
    }
}
\end{lstlisting}

\section{最小生成树}\label{ux6700ux5c0fux751fux6210ux6811}

\subsection{Kruskal}\label{kruskal}

按边权排序，用 DSU 合并不同连通块；恰好选 \passthrough{\lstinline!V-1!} 条边则图连通。复杂度 \passthrough{\lstinline!O(E log E)!}。可用于``最小化最大边''\,``第二小生成树''的基础。

\subsection{Prim}\label{prim}

从任意点出发，每次选未加入点中连接代价最小者，堆优化 \passthrough{\lstinline!O(E log V)!}；稠密图可用朴素 \passthrough{\lstinline!O(V²)!}。

MST 性质：若一条非树边加入形成环，环上最大边被替换可得到候选；用于瓶颈路径、边权敏感性、树上倍增求环上最大值。

\begin{lstlisting}[language={C++}]
sort(edges.begin(), edges.end(),
     [](auto &a, auto &b) { return a.w < b.w; });
DSU dsu(n); long long cost = 0; int used = 0;
for (auto e : edges) if (dsu.unite(e.u, e.v)) {
    cost += e.w; ++used;
}
bool connected = (used == n - 1);
\end{lstlisting}

\begin{lstlisting}[language={C++}]
vector<long long> best(n, INF); vector<char> used_prim(n);
priority_queue<pair<long long, int>, vector<pair<long long, int>>, greater<>> pq2;
best[0] = 0; pq2.push({0, 0}); long long mst = 0;
while (!pq2.empty()) {
    auto [w, u] = pq2.top(); pq2.pop();
    if (used_prim[u]) continue;
    used_prim[u] = true; mst += w;
    for (auto [v, c] : weighted_graph[u]) if (!used_prim[v] && c < best[v])
        best[v] = c, pq2.push({c, v});
}
\end{lstlisting}

\section{强连通分量（Strongly Connected Components，SCC）与缩点}\label{ux5f3aux8fdeux901aux5206ux91cfstrongly-connected-componentssccux4e0eux7f29ux70b9}

Tarjan：DFS 维护 \passthrough{\lstinline!dfn/low!} 与栈，\passthrough{\lstinline!low[u]==dfn[u]!} 时弹出一个 SCC，\passthrough{\lstinline!O(V+E)!}。Kosaraju：正图后序 + 反图 DFS，代码更直观。SCC 缩成 DAG 后做 DP、计数、可达性。缩点后不要把同一分量内的边当作 DAG 边。

\begin{lstlisting}[language={C++}]
int timer = 0; vector<int> dfn(n), low(n), stk, in(n), comp(n); int cc = 0;
void tarjan(int u) {
    dfn[u] = low[u] = ++timer; stk.push_back(u); in[u] = 1;
    for (int v : g[u]) {
        if (!dfn[v]) tarjan(v), low[u] = min(low[u], low[v]);
        else if (in[v]) low[u] = min(low[u], dfn[v]);
    }
    if (low[u] != dfn[u]) return;
    ++cc;
    while (true) {
        int v = stk.back(); stk.pop_back(); in[v] = 0; comp[v] = cc;
        if (v == u) break;
    }
}
\end{lstlisting}

\section{割点、桥、双连通}\label{ux5272ux70b9ux6865ux53ccux8fdeux901a}

无向图 Tarjan：\passthrough{\lstinline!low[v] >= dfn[u]!} 时，\passthrough{\lstinline!u!} 是割点分隔子树；\passthrough{\lstinline!low[v] > dfn[u]!} 时 \passthrough{\lstinline!(u,v)!} 是桥。根节点成为割点的条件是 DFS 子树数大于 1。边有重边时必须用边 id 判断父边，不能只判断父节点。

点双/边双连通可在 Tarjan 过程中用边栈或 DSU 缩桥。桥树上路径问题常转为树上 LCA/差分。

\begin{lstlisting}[language={C++}]
int timer = 0; vector<int> dfn(n), low(n); vector<pair<int, int>> bridges;
void dfs(int u, int parent_edge) {
    dfn[u] = low[u] = ++timer;
    for (auto [v, id] : g[u]) {
        if (id == parent_edge) continue;
        if (!dfn[v]) {
            dfs(v, id); low[u] = min(low[u], low[v]);
            if (low[v] > dfn[u]) bridges.push_back({u, v});
        } else low[u] = min(low[u], dfn[v]);
    }
}
\end{lstlisting}

\section{二分图}\label{ux4e8cux5206ux56fe}

\subsection{染色判定}\label{ux67d3ux8272ux5224ux5b9a}

BFS/DFS 交替染色；发现相邻同色则不是二分图。每个连通块可独立选择两种颜色，计数时乘法合并。

\begin{lstlisting}[language={C++}]
vector<int> color(n, -1);
bool ok = true;
for (int s = 0; s < n; ++s) if (color[s] < 0) {
    queue<int> q; q.push(s); color[s] = 0;
    while (!q.empty()) {
        int u = q.front(); q.pop();
        for (int v : g[u]) {
            if (color[v] < 0) color[v] = color[u] ^ 1, q.push(v);
            else if (color[v] == color[u]) ok = false;
        }
    }
}
\end{lstlisting}

\subsection{最大匹配}\label{ux6700ux5927ux5339ux914d}

\begin{itemize}
\tightlist
\item
  匈牙利 DFS：适合 \passthrough{\lstinline!V,E≤几千!}，复杂度 \passthrough{\lstinline!O(VE)!}。
\item
  Hopcroft--Karp：分层增广，\passthrough{\lstinline!O(E√V)!}，适合更大二分图。
\item
  二分图最小点覆盖 = 最大匹配（Kőnig）；最大独立集 = \passthrough{\lstinline!V - 最大匹配!}。
\end{itemize}

\begin{lstlisting}[language={C++}]
vector<int> matchR(n, -1), vis(n);
bool aug(int u) {
    for (int v : g[u]) if (vis[v] != stamp) {
        vis[v] = stamp;
        if (matchR[v] == -1 || aug(matchR[v]))
            return matchR[v] = u, true;
    }
    return false;
}
int matching = 0;
for (int u = 0; u < left_n; ++u) ++stamp, matching += aug(u);
\end{lstlisting}

\subsection{带权匹配}\label{ux5e26ux6743ux5339ux914d}

KM 算法求完备二分图最大权匹配，通常约 \passthrough{\lstinline!O(n³)!}；需要处理负权时整体平移或按模板维护顶标。

\section{网络流}\label{ux7f51ux7edcux6d41}

\subsection{建模}\label{ux5efaux6a21}

容量表示资源上限，源点连供应端，需求端连汇点。拆点处理点容量；上下界流增加下界需求并用超级源汇检查可行性。

\subsection{Dinic}\label{dinic}

BFS 建层次图，DFS 在层次递增的图上发送阻塞流；一般题目规模足够使用，二分图匹配也可套用。每条边要维护反向边 \passthrough{\lstinline!cap!}。

\begin{lstlisting}[language={C++}]
struct Edge { int to, rev; long long cap; };
vector<vector<Edge>> fg;
vector<int> level, it;
void add_edge(int u, int v, long long c) {
    Edge a{v, (int)fg[v].size(), c}, b{u, (int)fg[u].size(), 0};
    fg[u].push_back(a); fg[v].push_back(b);
}
bool bfs(int s, int t) {
    fill(level.begin(), level.end(), -1); queue<int> q;
    level[s] = 0; q.push(s);
    while (!q.empty()) { int u = q.front(); q.pop();
        for (auto &e : fg[u]) if (e.cap && level[e.to] < 0)
            level[e.to] = level[u] + 1, q.push(e.to);
    }
    return level[t] >= 0;
}
long long dfs(int u, int t, long long f) {
    if (u == t) return f;
    for (int &i = it[u]; i < (int)fg[u].size(); ++i) {
        Edge &e = fg[u][i];
        if (e.cap && level[e.to] == level[u] + 1) {
            long long got = dfs(e.to, t, min(f, e.cap));
            if (got) { e.cap -= got; fg[e.to][e.rev].cap += got; return got; }
        }
    }
    return 0;
}
\end{lstlisting}

最大流最小割：最大流值等于最小割容量；割边可由残量网络中从源可达的点集确定。最小割常用于``选/不选且有冲突代价''的闭合图建模。

\subsection{费用流}\label{ux8d39ux7528ux6d41}

增广最短路（势能 + Dijkstra 或 SPFA）反复找最小费用路径；适合规模较小的带费用匹配/运输，注意负边与势能初始化。

\begin{lstlisting}[language={C++}]
struct MCEdge { int to, rev, cap, cost; };
vector<vector<MCEdge>> mg(n);
void add_edge(int u, int v, int cap, int cost) {
    MCEdge a{v, (int)mg[v].size(), cap, cost};
    MCEdge b{u, (int)mg[u].size(), 0, -cost};
    mg[u].push_back(a); mg[v].push_back(b);
}
pair<int, long long> min_cost_flow(int s, int t, int need) {
    int flow = 0; long long cost = 0;
    vector<long long> pot(n), dist(n);
    vector<int> pv(n), pe(n);
    while (flow < need) {
        fill(dist.begin(), dist.end(), (long long)4e18);
        priority_queue<pair<long long, int>, vector<pair<long long, int>>,
                       greater<>> pq;
        dist[s] = 0; pq.push({0, s});
        while (!pq.empty()) {
            auto [du, u] = pq.top(); pq.pop();
            if (du != dist[u]) continue;
            for (int i = 0; i < (int)mg[u].size(); ++i) {
                auto const &e = mg[u][i];
                if (!e.cap) continue;
                long long nd = du + e.cost + pot[u] - pot[e.to];
                if (nd < dist[e.to])
                    dist[e.to] = nd, pv[e.to] = u, pe[e.to] = i,
                    pq.push({nd, e.to});
            }
        }
        if (dist[t] == (long long)4e18) break;
        for (int v = 0; v < n; ++v) if (dist[v] < (long long)4e18) pot[v] += dist[v];
        int add = need - flow;
        for (int v = t; v != s; v = pv[v]) add = min(add, mg[pv[v]][pe[v]].cap);
        for (int v = t; v != s; v = pv[v]) {
            auto &e = mg[pv[v]][pe[v]];
            e.cap -= add; mg[v][e.rev].cap += add; cost += 1LL * add * e.cost;
        }
        flow += add;
    }
    return {flow, cost};
}
\end{lstlisting}

\section{树上算法}\label{ux6811ux4e0aux7b97ux6cd5}

\subsection{DFS 序与子树}\label{dfs-ux5e8fux4e0eux5b50ux6811}

\passthrough{\lstinline!tin[u]!}、\passthrough{\lstinline!tout[u]!} 使子树变成连续区间 \passthrough{\lstinline![tin[u], tout[u]]!}。子树加/和直接用树状数组或线段树；路径则还需 LCA/重链。

\begin{lstlisting}[language={C++}]
int timer = 0; vector<int> tin(n), tout(n), depth(n), euler(n + 1);
void dfs(int u, int p) {
    tin[u] = ++timer; euler[timer] = u;
    for (int v : tree[u]) if (v != p)
        depth[v] = depth[u] + 1, dfs(v, u);
    tout[u] = timer;
}
\end{lstlisting}

\subsection{最近公共祖先（Lowest Common Ancestor，LCA）倍增}\label{ux6700ux8fd1ux516cux5171ux7956ux5148lowest-common-ancestorlcaux500dux589e}

预处理 \passthrough{\lstinline!up[u][j]!} 与深度，查询时先抬高深度差，再从大到小跳；预处理 \passthrough{\lstinline!O(V log V)!}，查询 \passthrough{\lstinline!O(log V)!}。同时可维护路径最小/最大/和。

\begin{lstlisting}[language={C++}]
const int LOG = 20;
int up[MAXN][LOG], dep[MAXN];
void init_lca(int u, int p) {
    up[u][0] = p;
    for (int j = 1; j < LOG; ++j) up[u][j] = up[up[u][j - 1]][j - 1];
    for (int v : tree[u]) if (v != p) dep[v] = dep[u] + 1, init_lca(v, u);
}
int lca(int u, int v) {
    if (dep[u] < dep[v]) swap(u, v);
    int d = dep[u] - dep[v];
    for (int j = 0; j < LOG; ++j) if (d >> j & 1) u = up[u][j];
    if (u == v) return u;
    for (int j = LOG - 1; j >= 0; --j)
        if (up[u][j] != up[v][j]) u = up[u][j], v = up[v][j];
    return up[u][0];
}
\end{lstlisting}

\subsection{树上差分}\label{ux6811ux4e0aux5deeux5206}

点路径 \passthrough{\lstinline!(u,v)!} 加一：\passthrough{\lstinline!d[u]++, d[v]++, d[lca]--, d[parent(lca)]--!}；边路径把 \passthrough{\lstinline!d[lca] -= 2!}。后序累加得到每点/边被经过次数。

\begin{lstlisting}[language={C++}]
void add_path(int u, int v) {
    int w = lca(u, v);
    ++diff[u]; ++diff[v]; --diff[w];
    if (up[w][0]) --diff[up[w][0]];
}
void collect(int u, int p) {
    for (int v : tree[u]) if (v != p)
        collect(v, u), diff[u] += diff[v];
}
\end{lstlisting}

\subsection{直径与重心}\label{ux76f4ux5f84ux4e0eux91cdux5fc3}

树直径：从任一点 BFS/DFS 找最远 \passthrough{\lstinline!s!}，再从 \passthrough{\lstinline!s!} 找最远 \passthrough{\lstinline!t!}；\passthrough{\lstinline!dist(s,t)!} 为直径。带权树同理（边权非负）。重心：删除后最大连通块最小的点，可由子树大小判断。

\begin{lstlisting}[language={C++}]
pair<int, long long> farthest(int s) {
    vector<long long> d(n, -1); queue<int> q;
    d[s] = 0; q.push(s);
    int best = s;
    while (!q.empty()) { int u = q.front(); q.pop();
        if (d[u] > d[best]) best = u;
        for (auto [v, w] : weighted_tree[u]) if (d[v] < 0)
            d[v] = d[u] + w, q.push(v);
    }
    return {best, d[best]};
}
\end{lstlisting}

\subsection{树形 DP}\label{ux6811ux5f62-dp}

先子树后父亲。常见状态：选/不选（独立集）、子树颜色、背包合并、最大匹配、距离和。换根 DP：第一次 DFS 求根答案，第二次用父节点答案推出子节点答案。

\begin{lstlisting}[language={C++}]
void tree_dp(int u, int p) {
    dp[u][0] = 0; dp[u][1] = weight[u];
    for (int v : tree[u]) if (v != p) {
        tree_dp(v, u);
        dp[u][0] += max(dp[v][0], dp[v][1]);
        dp[u][1] += dp[v][0];
    }
}
\end{lstlisting}

\subsection{重链剖分（Heavy-Light Decomposition，HLD）}\label{ux91cdux94feux5256ux5206heavy-light-decompositionhld}

把每条重链映射到数组，路径拆成 \passthrough{\lstinline!O(log V)!} 个连续区间，配合线段树实现路径/子树修改查询，复杂度 \passthrough{\lstinline!O(log² V)!}。先求 \passthrough{\lstinline!size, heavy!}，再求 \passthrough{\lstinline!top, dfn!}。

\begin{lstlisting}[language={C++}]
void dfs1(int u, int p) {
    parent[u] = p; size[u] = 1; heavy[u] = 0;
    for (int v : tree[u]) if (v != p) {
        dep[v] = dep[u] + 1; dfs1(v, u); size[u] += size[v];
        if (!heavy[u] || size[v] > size[heavy[u]]) heavy[u] = v;
    }
}
void dfs2(int u, int topf) {
    top[u] = topf; dfn[u] = ++timer;
    if (heavy[u]) dfs2(heavy[u], topf);
    for (int v : tree[u]) if (v != parent[u] && v != heavy[u]) dfs2(v, v);
}
long long path_sum(int u, int v) {
    long long ans = 0;
    while (top[u] != top[v]) {
        if (dep[top[u]] < dep[top[v]]) swap(u, v);
        ans += seg.query(dfn[top[u]], dfn[u]); u = parent[top[u]];
    }
    if (dep[u] > dep[v]) swap(u, v);
    return ans + seg.query(dfn[u], dfn[v]);
}
\end{lstlisting}

\subsection{点分治（入门）}\label{ux70b9ux5206ux6cbbux5165ux95e8}

以重心分治树，统计跨重心的路径信息，递归处理子树；每层总计 \passthrough{\lstinline!O(V)!}，深度 \passthrough{\lstinline!O(log V)!}，常见复杂度 \passthrough{\lstinline!O(V log V)!}。适用于树上距离小于/等于 k 的点对计数、距离最小值。

\begin{lstlisting}[language={C++}]
int get_size(int u, int p) {
    sz[u] = 1;
    for (int v : tree[u]) if (v != p && !removed[v]) sz[u] += get_size(v, u);
    return sz[u];
}
int get_centroid(int u, int p, int total) {
    for (int v : tree[u]) if (v != p && !removed[v])
        if (sz[v] * 2 > total) return get_centroid(v, u, total);
    return u;
}
void divide(int entry) {
    int c = get_centroid(entry, 0, get_size(entry, 0));
    removed[c] = true;
    // 统计经过 c 的贡献
    for (int v : tree[c]) if (!removed[v]) divide(v);
}
\end{lstlisting}

\subsection{树上启发式合并（Disjoint Set Union on Tree，DSU on Tree）}\label{ux6811ux4e0aux542fux53d1ux5f0fux5408ux5e76disjoint-set-union-on-treedsu-on-tree}

先求子树大小和重儿子；轻儿子递归后清空，重儿子保留统计结果，再把所有轻子树合并进来。每个节点被清除/加入的次数受重轻性质控制，总复杂度通常为 \passthrough{\lstinline!O(V log V)!}。

\begin{lstlisting}[language={C++}]
vector<int> sz(n + 1), heavy(n + 1); // 子树大小、重儿子。
vector<int> color(n + 1), freq(MAXC); // 颜色和当前统计集合中的频次。
vector<int> answer(n + 1);
int distinct = 0; // 示例维护量：当前集合中出现过的颜色数。

void dfs_size(int u, int p) {
    // 第一次 DFS 只决定轻/重儿子，不修改答案统计。
    sz[u] = 1;
    for (int v : tree[u]) if (v != p) {
        dfs_size(v, u); sz[u] += sz[v];
        if (!heavy[u] || sz[v] > sz[heavy[u]]) heavy[u] = v;
    }
}
void add_subtree(int u, int p, int delta) {
    // delta=+1 是加入，delta=-1 是清空。
    // 这里的“答案维护器”可以替换成最大频率、频率平方和等任何可增删统计。
    if (delta == 1 && freq[color[u]]++ == 0) ++distinct;
    if (delta == -1 && --freq[color[u]] == 0) --distinct;
    // 必须递归整棵子树；只处理 u 会漏掉大量颜色。
    for (int v : tree[u]) if (v != p) add_subtree(v, u, delta);
}
void dsu_on_tree(int u, int p, bool keep) {
    // 轻儿子先算并清空：它们只贡献一次，避免统计互相污染。
    for (int v : tree[u]) if (v != p && v != heavy[u]) dsu_on_tree(v, u, false);
    // 重儿子保留答案，后面把轻子树和 u 合并进这份大集合。
    if (heavy[u]) dsu_on_tree(heavy[u], u, true);
    // 把所有轻儿子整棵加入重儿子的统计集合。
    for (int v : tree[u]) if (v != p && v != heavy[u]) add_subtree(v, u, 1);
    // 最后加入 u 自己，此时集合恰好是 u 的整棵子树。
    if (freq[color[u]]++ == 0) ++distinct;
    answer[u] = distinct; // 换成题目所需的当前子树统计量
    if (!keep) {
        // u 是父节点的轻儿子时，当前统计不会被父节点复用，整棵删除。
        add_subtree(u, p, -1);
    }
}
\end{lstlisting}

\subsection{长链剖分}\label{ux957fux94feux5256ux5206}

长链剖分按``向下最长深度''而不是子树大小选重儿子，适合树上深度/祖先方向的 DP 优化和长链合并。它和 HLD 的选重标准不同，不能混用复杂度证明。

\begin{lstlisting}[language={C++}]
vector<int> parent2(n + 1), depth2(n + 1), heavy_len(n + 1);
vector<int> heavy2(n + 1), top2(n + 1);
void long_chain_dfs1(int u, int p) {
    // heavy_len[u] 是从 u 向下的最长链长度，不是子树大小。
    parent2[u] = p; heavy_len[u] = 1;
    for (int v : tree[u]) if (v != p) {
        depth2[v] = depth2[u] + 1;
        long_chain_dfs1(v, u);
        // 选择向下最深的儿子作为长链儿子。
        if (!heavy2[u] || heavy_len[v] > heavy_len[heavy2[u]]) heavy2[u] = v;
    }
    if (heavy2[u]) heavy_len[u] = heavy_len[heavy2[u]] + 1;
}
void long_chain_dfs2(int u, int topf) {
    // 长儿子沿用同一个链顶，其他儿子各自开新链。
    top2[u] = topf;
    if (heavy2[u]) long_chain_dfs2(heavy2[u], topf);
    for (int v : tree[u])
        if (v != parent2[u] && v != heavy2[u]) long_chain_dfs2(v, v);
}
\end{lstlisting}

\subsection{Link-Cut Tree 动态树（LCT）}\label{link-cut-tree-ux52a8ux6001ux6811lct}

LCT 用 splay 维护动态森林，支持 \passthrough{\lstinline!link/cut!}、换根和路径聚合。下面模板维护路径和；将 \passthrough{\lstinline!pull!} 中的 \passthrough{\lstinline!sum!} 换成结构体即可维护最大值、异或等可合并信息。

\begin{lstlisting}[language={C++}]
struct LCT {
    struct Node {
        int ch[2]{};       // splay 左右儿子。
        int fa{};          // 辅助树父亲；不等同于原树父亲。
        long long val{};   // 点权。
        long long sum{};   // 当前 splay 子树的维护信息。
        bool rev{};        // 是否有待下传的换根翻转标记。
    };
    vector<Node> t;
    LCT(int n): t(n + 1) {}
    bool is_root(int x) const {
        // x 没有 splay 父亲，或不是父亲的直接儿子时，是辅助树根。
        int f = t[x].fa;
        return !f || (t[f].ch[0] != x && t[f].ch[1] != x);
    }
    void pull_sum(int x) {
        // 修改儿子后必须重新计算聚合信息。
        t[x].sum = t[x].val + t[t[x].ch[0]].sum + t[t[x].ch[1]].sum;
    }
    void apply_rev(int x) {
        if (!x) return;
        // 翻转只交换辅助树儿子；原树路径方向由 makeroot 实现。
        swap(t[x].ch[0], t[x].ch[1]); t[x].rev ^= 1;
    }
    void push(int x) {
        // 旋转前必须把从辅助树根到 x 的所有 rev 标记按顺序下传。
        if (t[x].rev) apply_rev(t[x].ch[0]), apply_rev(t[x].ch[1]), t[x].rev = false;
    }
    void rotate(int x) {
        // 一次旋转把 x 提升；b 是 x 原来靠旋转方向的另一棵子树。
        int y = t[x].fa, z = t[y].fa, dx = (t[y].ch[1] == x), b = t[x].ch[dx ^ 1];
        if (!is_root(y)) t[z].ch[t[z].ch[1] == y] = x;
        t[x].fa = z; t[x].ch[dx ^ 1] = y; t[y].fa = x; t[y].ch[dx] = b;
        if (b) t[b].fa = y;
        pull_sum(y); pull_sum(x);
    }
    void splay(int x) {
        // 先收集祖先，再从上到下 push，确保旋转时方向标记正确。
        vector<int> st{ x };
        for (int y = x; !is_root(y); y = t[y].fa) st.push_back(t[y].fa);
        while (!st.empty()) push(st.back()), st.pop_back();
        while (!is_root(x)) {
            int y = t[x].fa, z = t[y].fa;
            if (!is_root(y)) ((t[y].ch[0] == x) ^ (t[z].ch[0] == y)) ? rotate(x) : rotate(y);
            rotate(x);
        }
    }
    void access(int x) {
        // 依次把 x 到原树根的 preferred path 变成一条辅助树路径。
        // y 是上一轮处理的节点，挂到 x 的右儿子。
        for (int y = 0; x; x = t[y = x].fa) splay(x), t[x].ch[1] = y, pull_sum(x);
    }
    void makeroot(int x) {
        // access+splay 后 x 表示整条根到 x 的路径，翻转它即可让 x 成为原树根。
        access(x); splay(x); apply_rev(x);
    }
    int findroot(int x) {
        // 暴露 x 后不断向左走，左端点就是所在原树的根。
        access(x); splay(x);
        while (t[x].ch[0]) push(x), x = t[x].ch[0];
        splay(x); return x;
    }
    void split(int x, int y) {
        // 之后 y 的 splay 子树正好维护原树路径 x..y。
        makeroot(x); access(y); splay(y);
    }
    void link(int x, int y) {
        // 先换根保证 x 没有原树父亲，再确认两点不连通。
        makeroot(x); if (findroot(y) != x) t[x].fa = y;
    }
    void cut(int x, int y) {
        // 暴露边 x-y；合法形态要求 y 的左儿子是 x 且 x 没有右儿子。
        split(x, y);
        if (t[y].ch[0] == x && !t[x].ch[1]) t[y].ch[0] = t[x].fa = 0, pull_sum(y);
    }
    long long path_sum(int x, int y) {
        split(x, y);
        return t[y].sum;
    }
};
\end{lstlisting}

\section{欧拉路、欧拉回路与哈密顿误区}\label{ux6b27ux62c9ux8defux6b27ux62c9ux56deux8defux4e0eux54c8ux5bc6ux987fux8befux533a}

无向图存在欧拉回路当且仅当非零度点连通且所有度为偶；欧拉路径有且仅有 0 或 2 个奇度点。有向图对应入度=出度，路径起点出度比入度多 1、终点相反。Hierholzer 用栈在线性时间构造。哈密顿路径一般是 NP-hard，不要把两者混淆。

\begin{lstlisting}[language={C++}]
vector<int> ptr(n), euler;
void hierholzer(int u) {
    while (ptr[u] < (int)g[u].size()) {
        auto [v, id] = g[u][ptr[u]++];
        if (used[id]) continue;
        used[id] = true; hierholzer(v);
    }
    euler.push_back(u);
}
\end{lstlisting}

\section{图论组合套路}\label{ux56feux8bbaux7ec4ux5408ux5957ux8def}

\begin{itemize}
\tightlist
\item
  ``删掉一些边使图满足条件''常考虑 MST、割、桥、最小割。
\item
  ``两点之间经过某类边最少''常做分层图/状态图。
\item
  ``最多经过 k 次特殊边''建 \passthrough{\lstinline!k+1!} 层节点。
\item
  ``必须恰好经过/使用''常用流、DP 或矩阵，不要只跑最短路。
\item
  DAG + 计数要同时维护 \passthrough{\lstinline!ways[v]!}，相同最优值时累加、更新更优时覆盖。
\end{itemize}

